\section{Protocolo experimental y reproducibilidad}
\subsection{Familias y condiciones}
Se distinguen \textbf{A3}: CNN inicial y experimentos aleatorios/DM;
\textbf{A4}: MLP, CNN original, CNN ampliada y CNN AdamW/LR; y
\textbf{F}: hiperparámetros históricos reutilizados en el A3 corregido y
conservados en el reporte final. Las corridas A4 y el par F aquí analizado
declaran seed 0, train completo, la misma validación, entrada 112,
batch 64 y hasta 150 épocas. No se promedian resultados entre familias.

En A4, las CNN original y ampliada usan Adam LR $10^{-3}$ y paciencia de cinco
evaluaciones sin mejora de F1. La MLP usa Adam LR $10^{-4}$, paciencia diez y
mínimo diez épocas, después de un ensayo fallido con $10^{-3}$. La CNN AdamW/LR
usa weight decay 0.01 sin sesgos, reducción por pérdida de validación de factor
0.5, paciencia dos, mínimo $10^{-5}$ y parada temprana de doce evaluaciones.
Todos los A4 usan dropout 0.5 y aumentos: su comparación Adam/AdamW no equivale
a retirar toda la regularización explícita.

En F se compara dropout 0.5, aumentos y AdamW/WD 0.01 frente a dropout 0,
sin aumentos y Adam/WD 0. Ambas conservan LR inicial $10^{-3}$, la misma regla
de reducción de tasa, paciencia doce y tope 150. Sus trayectorias reales de LR
pueden diferir. El checkpoint se selecciona por máximo F1 macro; exactitud y
pérdida se asocian a ese mismo modelo, no a máximos independientes.
El mínimo de pérdida de todo un historial es otra cantidad, obtenida posiblemente
en otra época; no se usa para completar métricas del checkpoint por F1.

\subsection{Medición de la brecha}
En A4, train se mide durante actualizaciones con dropout y aumentos. En F se
registran además las primeras 3,712 imágenes fijas de train, sin aumentos y en
modo inferencia al final de cada época. Solo estas métricas limpias permiten
comparar con validación bajo el mismo modo de evaluación:
\begin{equation}
\Delta F=F1_{\rm train,limpio}-F1_{\rm val}.
\end{equation}
Ese subconjunto no es todo train ni una muestra estratificada. Sus métricas
se recuperan de históricos, sin predicciones individuales de train para repetir
el recálculo. Se comparan mejores checkpoints y, como complemento, la época
34 registrada en ambas corridas F. No se interpolan valores ausentes.

\subsection{Artefactos y límites de ejecución}
Los originales A4 están en
\path{reports/entregable_4/avance4_2026-10-21}; el par F, en
\path{reports/final/experimentos/runs}. Los derivados registran IDs de familia,
rutas y hashes en carpetas de corrección del A4. La tabla completa por clase
se obtiene de predicciones, sin modificar sus fuentes.

El notebook \path{notebooks/equipo_proyecto_avance4.ipynb} del
\href{https://github.com/Lunaaaalj/NeuralNetworksTeamEnvironment}{repositorio del equipo}
separa análisis histórico de entrenamiento. El primer modo usa CSV/JSON sin
frameworks de entrenamiento ni TFRecords; el segundo exige datos y dependencias
y escribe en una salida nueva. En el entorno de revisión faltaron TFRecords y
frameworks, por lo que no se reentrenaron modelos ni se recalcularon pérdidas
sobre imágenes. Una semilla no estima variabilidad; validación también se usa
para seleccionar modelos y no constituye test independiente.
